\section{Finite-dimensional weights and sparse block operators}
\label{sec:weights}

The sparse assignment from \Cref{sec:sparse} will now produce the three
bounds required by \Cref{prop:direct-sum-obstruction}: a bound on each
block row of \(T\), each block column of \(A-T\), and every conjugation
difference \(T-U(g)TU(g)^{-1}\). All three bounds will be independent
of the number of labels. To make the energy grow faster than the fibre
dimension, we use rank-one projections onto many unit vectors in a
smaller space.

We first construct vectors whose synthesis operators are bounded on
every small subfamily. The proof uses the standard concentration,
finite-net and union-bound method for random matrices; compare
\cite[Section~5]{BDDW2008}. Here only an upper bound is needed, and the
argument below supplies its constants directly.

\begin{lemma}[An upper bound for all small subfamilies]
\label{lem:frames}
Let $n\ge 2$ and $r\ge 1$ be integers, and put
\begin{equation}
\label{eq:frame-parameters}
 p=\min(2r,n),\qquad
 k=\left\lceil 2r\log(2n)\right\rceil,
 \qquad b_0=10,
\end{equation}
where $\log$ denotes the natural logarithm.  There are unit vectors
$v_1,\ldots,v_n\in\C^k$ such that, for every
$J\subseteq\{1,\ldots,n\}$ with $|J|\le p$, the operator
\[
 V_J\colon\ell^2(J)\longrightarrow\C^k,
 \qquad V_J z=\sum_{i\in J}z_i v_i,
\]
satisfies $\norm{V_J}\le b_0$.
\end{lemma}

\begin{proof}
Choose independent random variables $\varepsilon_{a,i}$, for
$1\le a\le k$ and $1\le i\le n$, each taking the values $1$ and $-1$
with probability $1/2$, and set
\[
 (v_i)_a=k^{-1/2}\varepsilon_{a,i}.
\]
Every $v_i$ has norm exactly one.  Fix a set
$J=\{i_1<\cdots<i_p\}$ of cardinality $p$, and fix real unit vectors
$w\in\mathbb R^k$ and $z\in\mathbb R^p$.  The random variable
\[
 X=w^{\mathsf T}V_Jz
   =\frac1{\sqrt{k}}
     \sum_{a=1}^k\sum_{j=1}^p w_a z_j\varepsilon_{a,i_j}
\]
is a sum of independent signs with deterministic coefficients whose
squares sum to
\[
 \frac1k\sum_{a=1}^k\sum_{j=1}^p w_a^2z_j^2=\frac1k.
\]
The inequality $\cosh t\le\exp(t^2/2)$ consequently gives, for every
$t\in\mathbb R$,
\[
 \mathbb E\exp(tX)
 =\prod_{a,j}\cosh\!\left(\frac{t w_a z_j}{\sqrt{k}}\right)
 \le\exp\!\left(\frac{t^2}{2k}\right).
\]
For $u>0$, applying Markov's inequality to $\exp(tX)$ and choosing
$t=ku$ yields $\Pr(X>u)\le\exp(-ku^2/2)$.  The same argument for
$-X$ therefore proves
\begin{equation}
\label{eq:sign-tail}
 \Pr(|X|>u)\le 2\exp(-ku^2/2).
\end{equation}
In particular, the probability that $|X|>5$ is at most
$2\exp(-25k/2)$.

For each positive integer $m$, the real Euclidean unit sphere has a
$1/4$-net with at most $9^m$ points.  Indeed, take a maximal
$1/4$-separated set on that sphere.  The open balls of radius $1/8$
about its points are pairwise disjoint and lie in the ball of radius
$9/8$ about the origin.  Comparing their $m$-dimensional volumes bounds
the number of points by $9^m$.  Maximality makes the set a $1/4$-net.

Fix such nets $\mathcal N_k\subset\mathbb R^k$ and
$\mathcal N_p\subset\mathbb R^p$.  For any real $k\times p$ matrix
$M$, let $L$ be its operator norm on real Euclidean spaces.  If
$|w_0^{\mathsf T}Mz_0|\le5$ for all
$(w_0,z_0)\in\mathcal N_k\times\mathcal N_p$, then, for arbitrary
real unit vectors $w,z$, choose net points satisfying
$\norm{w-w_0}\le1/4$ and $\norm{z-z_0}\le1/4$.  We obtain
\[
 \left|w^{\mathsf T}Mz-w_0^{\mathsf T}Mz_0\right|
 \le \norm{w-w_0}L\norm{z}
       +\norm{w_0}L\norm{z-z_0}
 \le\frac L2.
\]
Taking the supremum over $w,z$ gives $L\le5+L/2$, hence $L\le10$.

There are at most $n^p$ sets $J$ of size $p$.  By
\Cref{eq:sign-tail} and the union bound, the probability that one of
the corresponding net inequalities fails is at most
\begin{equation}
\label{eq:frame-failure-probability}
 2n^p9^{k+p}\exp(-25k/2).
\end{equation}
The parameters satisfy
\[
 k\ge2r\log(2n)\ge p,
 \qquad k\ge p\log n.
\]
Here $\log(2n)\ge\log4>1$ and $p\le2r$.  Thus the logarithm of
the bound in \Cref{eq:frame-failure-probability} is at most
\[
 \log2+k+2k\log9-\frac{25k}{2}
 =\log2-\left(\frac{25}{2}-1-2\log9\right)k<0.
\]
For example, the coefficient in parentheses exceeds $7$, whereas
$k\ge1$ and $\log2<1$.  The bound in
\Cref{eq:frame-failure-probability} is therefore strictly less than
one.  There is a realization for which every size-$p$ synthesis
matrix has real operator norm at most $10$.

For a real matrix $M$ with real operator norm $L$, and real vectors
$u,v$, one has
\[
 \norm{M(u+iv)}^2=\norm{Mu}^2+\norm{Mv}^2
 \le L^2\bigl(\norm{u}^2+\norm{v}^2\bigr).
\]
Since real vectors are also complex vectors, this proves equality of
the real and complex operator norms.  The realization just obtained
therefore has the claimed complex operator norm bounds.  Finally,
every set with fewer than $p$ indices is contained in one with $p$
indices, and its synthesis operator is the restriction to a coordinate
subspace of that larger synthesis operator.  This proves the assertion
for all $|J|\le p$.
\end{proof}

The vectors just constructed control each small set of incident labels.
We next transfer that finite-dimensional estimate to an operator on the
group. The intermediate space keeps one scalar coordinate for each
labelled edge, even when several labels describe the same group element.

\subsection{A factorization through the indexed edges}

Let $G$ be a countable group, let $s_1,\ldots,s_n\in G$ be a finite
list, and allow the same group element to occur at several distinct
indices.  Let $v_1,\ldots,v_n\in\C^k$ be unit vectors and put
\[
 P_i=v_i v_i^*,\qquad
 \mathcal H=\ell^2(G;\C^k).
\]
Thus $P_i$ is a rank-one orthogonal projection and
$\norm{P_i}=1$.  For a function
$c\colon G\times\{1,\ldots,n\}\to\C$ with $|c(x,i)|\le1$, define
\begin{equation}
\label{eq:masked-operator}
 (T_c\xi)(x)=\sum_{i=1}^n c(x,i)P_i\xi(xs_i),
 \qquad \xi\in\mathcal H.
\end{equation}
For each fixed $i$, the corresponding summand is the composition of
the unitary right shift $\xi(x)\mapsto\xi(xs_i)$, the pointwise
application of $P_i$, and multiplication by the scalar function
$c(\cdot,i)$.  Its norm is at most one.  In particular,
\Cref{eq:masked-operator} defines a bounded operator and
$\norm{T_c}\le n$, without any sparsity assumption.  The next lemma
gives a bound independent of $n$ under two sparsity assumptions.

\begin{lemma}[The masked operator estimate]
\label{lem:masked}
In the preceding setting, let $1\le p\le n$ and $b_0>0$.  Suppose
that
\[
 \norm{V_J}\le b_0
 \quad\text{for every }J\subseteq\{1,\ldots,n\}
 \text{ with }|J|\le p,
\]
where $V_J$ is the synthesis operator with columns $v_i$, $i\in J$.
For a mask $c$ as in \Cref{eq:masked-operator}, define
\[
 \mathcal R_x=\{i:c(x,i)\ne0\},\qquad
 \mathcal C_y=\{i:c(ys_i^{-1},i)\ne0\}.
\]
If $|\mathcal R_x|\le p$ for every $x\in G$ and $|\mathcal C_y|\le p$ for every
$y\in G$, then
\[
 \norm{T_c}\le b_0^2.
\]
\end{lemma}

\begin{proof}
Keep the index $i$ as part of each edge and set
\[
 E_c=\{(x,i):c(x,i)\ne0\},
 \qquad \mathcal E_c=\ell^2(E_c).
\]
Define the analysis operator $F\colon\mathcal H\to\mathcal E_c$,
the diagonal multiplication operator
$M_c\colon\mathcal E_c\to\mathcal E_c$, and the synthesis operator
$S\colon\mathcal E_c\to\mathcal H$ by
\begin{align*}
 (F\xi)(x,i)&=v_i^*\xi(xs_i),\\
 (M_c\alpha)(x,i)&=c(x,i)\alpha(x,i),\\
 (S\alpha)(x)&=\sum_{i\in \mathcal R_x}v_i\alpha(x,i).
\end{align*}
Figure~\ref{fig:edge-factorization} illustrates the role of the edge
coordinates. The following estimates both justify these operators on the
stated Hilbert spaces and bound their norms.
\begin{figure}[tb]
\centering
\begin{tikzpicture}[
  >={Stealth[length=2mm]},
  fibre/.style={draw,rounded corners,minimum width=23mm,minimum height=13mm,align=center},
  edgecoord/.style={draw,minimum width=27mm,minimum height=9mm,align=center},
  every node/.style={font=\small}]
\node[fibre] (column) at (0,0) {column fibre\\$\xi(y)\in\C^k$};
\node[edgecoord] (ei) at (4,1) {$v_i^*\xi(y)\in\C$};
\node[edgecoord] (ej) at (4,-1) {$v_j^*\xi(y)\in\C$};
\node[fibre] (row) at (8,0) {row fibre\\$(T_c\xi)(x)\in\C^k$};
\draw[->] (column) -- node[above,sloped] {$F$: label $i$} (ei);
\draw[->] (column) -- node[below,sloped] {$F$: label $j$} (ej);
\draw[->] (ei) -- node[above,sloped] {$c(x,i)v_i$} (row);
\draw[->] (ej) -- node[below,sloped] {$c(x,j)v_j$} (row);
\node at (4,1.75) {distinct edge coordinates};
\node at (4,-1.75) {$s_i=s_j$,\quad $i\ne j$};
\end{tikzpicture}

\caption{Two different labels with \(s_i=s_j\) give parallel edges from
row \(x\) to column \(y=xs_i=xs_j\). The operator acts from the column
fibre toward the row fibre. Analysis produces separate scalar
coordinates \((x,i)\) and \((x,j)\); multiplication by the mask and
synthesis then add their contributions in the row fibre. The picture
shows only these two coordinates of the factorization \(T_c=SM_cF\).}
\label{fig:edge-factorization}
\end{figure}  For $S$, the row assumption gives
\[
 \norm{S\alpha}^2
 =\sum_x\norm{V_{\mathcal R_x}(\alpha(x,i))_{i\in \mathcal R_x}}^2
 \le b_0^2\sum_x\sum_{i\in \mathcal R_x}|\alpha(x,i)|^2
 =b_0^2\norm{\alpha}^2.
\]
For $F$, regroup the nonnegative summands by the column $y=xs_i$.
For fixed $y$ and fixed label $i$, there is exactly one possible
source, namely $x=ys_i^{-1}$.  The edges arriving at $y$ are
therefore indexed by the set $\mathcal C_y$, with no repeated index, and
\begin{align*}
 \norm{F\xi}^2
 &=\sum_y\sum_{i\in \mathcal C_y}|v_i^*\xi(y)|^2\\
 &=\sum_y\norm{V_{\mathcal C_y}^*\xi(y)}^2
 \le b_0^2\sum_y\norm{\xi(y)}^2
 =b_0^2\norm{\xi}^2.
\end{align*}
Empty index sets contribute zero in both calculations.  These
estimates give $\norm{S}\le b_0$ and $\norm{F}\le b_0$, and
$|c(x,i)|\le1$ gives $\norm{M_c}\le1$.  Substitution into the
definitions yields the exact factorization
\[
 T_c=S M_c F.
\]
Consequently $\norm{T_c}\le b_0^2$.

The uniqueness of $x=ys_i^{-1}$ uses a fixed label $i$, not the
distinctness of the elements in the list $s_1,\ldots,s_n$.
In particular, if $s_i=s_j$ for $i\ne j$, the two indexed edges remain
distinct coordinates of $\mathcal E_c$.  Their contributions to a
common operator block are summed by the same factorization.  No
distinctness assumption on the products is needed.
\end{proof}

\subsection{Consequences of an assignment to endpoints}

We now apply the factorization three times to obtain the hypotheses of
\Cref{prop:direct-sum-obstruction}. An endpoint assignment is sparse
only for the ones at rows and the zeros at columns. Restricting to a
single row or column, and then subtracting a translated assignment,
will give masks that are sparse on both sides.

Fix integers $n\ge2$ and $r\ge1$, and choose $k,p,b_0$ and
$v_1,\ldots,v_n$ as in \Cref{lem:frames}.  Keep the preceding
notation $G,s_i,P_i,\mathcal H,T_c$.  Suppose a function
$a\colon G\times\{1,\ldots,n\}\to\{0,1\}$ satisfies
\begin{equation}
\label{eq:weights-orientation-bounds}
 \sum_{i=1}^n a(x,i)\le r\quad(x\in G),
 \qquad
 \sum_{i=1}^n\bigl(1-a(ys_i^{-1},i)\bigr)\le r\quad(y\in G).
\end{equation}
Thus the ones are sparse at rows and the zeros are sparse at columns,
with both counts taken over indexed edges.  Set
\begin{equation}
\label{eq:weights-A-T-U}
 A=T_{\one},\qquad T=T_a,\qquad
 (U(g)\xi)(x)=\xi(g^{-1}x)\quad(g,x\in G),
\end{equation}
where $\one(x,i)=1$.  Both $A$ and $T$ are bounded operators, each
with norm at most $n$. For $y\in G$, let
$\iota_y:\C^k\to\mathcal H$ be the coordinate inclusion, and write
$q_y=\iota_y^*$ for coordinate evaluation.

\begin{corollary}[Individual block rows and columns]
\label{cor:block-rows-columns}
Under \Cref{eq:weights-orientation-bounds}, one has
\[
 \norm{q_xT}\le b_0^2\quad(x\in G),
 \qquad
 \norm{(A-T)\iota_y}\le b_0^2\quad(y\in G).
\]
\end{corollary}

\begin{proof}
Write $\one_E$ for the indicator of a condition $E$.
For fixed $x_0,y_0\in G$, define the masks
\[
 c_{\mathrm r}(x,i)=\one_{\{x=x_0\}}a(x,i),\qquad
 c_{\mathrm c}(x,i)=\one_{\{xs_i=y_0\}}(1-a(x,i)).
\]
The first inequality in \Cref{eq:weights-orientation-bounds} gives
at most $r$ active edges in total for $c_{\mathrm r}$; the second
gives the same bound for $c_{\mathrm c}$.
Thus every row and every column of either mask has at most
$\min(r,n)\le p$ active indices, and \Cref{lem:masked} bounds both
associated operator norms by $b_0^2$. Their operators are
\[
 T_{c_{\mathrm r}}=\iota_{x_0}q_{x_0}T,\qquad
 T_{c_{\mathrm c}}=(A-T)\iota_{y_0}q_{y_0}.
\]
The first norm equals $\norm{q_{x_0}T}$ because $\iota_{x_0}$ is
isometric. The second equals $\norm{(A-T)\iota_{y_0}}$ because
$\norm{q_{y_0}}=1$ and $q_{y_0}\iota_{y_0}=I_{\C^k}$.
These are the asserted bounds.
\end{proof}

\begin{proposition}[Uniformly bounded conjugation differences]
\label{prop:uniform-differences}
Under \Cref{eq:weights-orientation-bounds}, the maps $U(g)$ in
\Cref{eq:weights-A-T-U} form a unitary representation of $G$,
and $A$ commutes with every $U(g)$.  Define
\[
 D(g)=T-U(g)TU(g)^{-1}\quad(g\in G).
\]
Then
\begin{equation}
\label{eq:uniform-difference-bound}
 \sup_{g\in G}\norm{D(g)}\le b_0^2=100.
\end{equation}
Moreover, $D(e)=0$ and
\begin{equation}
\label{eq:weights-cocycle}
 D(gh)=D(g)+U(g)D(h)U(g)^{-1}\quad(g,h\in G).
\end{equation}
\end{proposition}

\begin{proof}
Left multiplication by a group element is a bijection of $G$, so
each $U(g)$ is unitary with inverse $U(g^{-1})$, and direct
substitution gives $U(g)U(h)=U(gh)$.
Left and right translations
commute and $U(g)$ acts trivially on the fibre $\C^k$; alternatively,
direct substitution in \Cref{eq:masked-operator} with $c=\one$
gives $U(g)A=AU(g)$.

For a general mask $a$, direct calculation gives
\[
 (U(g)TU(g)^{-1}\xi)(x)
 =\sum_{i=1}^n a(g^{-1}x,i)P_i\xi(xs_i).
\]
Thus $D(g)=T_{c_g}$ for the mask
\[
 c_g(x,i)=a(x,i)-a(g^{-1}x,i)\in\{-1,0,1\}.
\]
At any row $x$, the active indices are contained in
\[
 \{i:a(x,i)=1\}\ \cup\ \{i:a(g^{-1}x,i)=1\}.
\]
Each set in this union has at most $r$ elements by
\Cref{eq:weights-orientation-bounds}.  Hence there are at most
$\min(2r,n)=p$ active indices at that row.

For the column estimate, write
\[
 Z_y=\{i:a(ys_i^{-1},i)=0\}\quad(y\in G).
\]
The second inequality in \Cref{eq:weights-orientation-bounds}
gives $|Z_y|\le r$.  The active indices of $c_g$ at a column $y$
are precisely those for which
\[
 a(ys_i^{-1},i)\ne a(g^{-1}ys_i^{-1},i).
\]
Since $a$ takes only the values zero and one, this set is
$Z_y\mathbin\triangle Z_{g^{-1}y}$, the symmetric difference of
the two zero sets.  It has at most $2r$ elements, and at most $n$,
so its cardinality is at most $p$.  As in \Cref{lem:masked},
fixed $y$ and fixed label $i$ specify the unique source
$ys_i^{-1}$.  These counts therefore concern exactly the indexed
edges of the mask, including when products coincide.

Both hypotheses of \Cref{lem:masked} now hold for $c_g$, and
$|c_g(x,i)|\le1$.  That lemma yields
$\norm{D(g)}\le b_0^2$ for every $g$, proving
\Cref{eq:uniform-difference-bound}.

Finally, $D(e)=T-T=0$. The cocycle identity in
\Cref{eq:weights-cocycle} follows from the calculation in the proof
of \Cref{prop:direct-sum-obstruction}.
\end{proof}
